% SPDX-License-Identifier: CC-BY-SA-4.0
% Copyright (c) 2026 Rafael Beraldo
% Parte 4: o radioamadorismo, sua história e suas entidades
% Incluído por presentation.tex; compile a partir dele.

% =============================================================================
\section{O radioamadorismo}
% =============================================================================

\begin{frame}{Radioamadorismo: hobby, ciência e comunidade}
  \begin{dica}
    O radioamadorismo é um \alert<1>{hobby científico}, \alert<2>{social} (“primeira rede social”) e \alert<3>{competitivo}, que tem papel importante no \alert<4>{desenvolvimento das telecomunicações} e, eventualmente, na \alert<4>{segurança civil}. \alert<5>{Radioamadores} são pessoas \alert<5>{licenciadas para usar e experimentar} com certas partes do espectro eletromagnético.

    \raggedleft
    \small
    \smallskip
    Com base em \citet{iaru2020}
  \end{dica}
\end{frame}

\begin{frame}{As origens do radioamadorismo}
  \begin{columns}
    \column{.45\textwidth}
    \begin{itemize}[<+->]
    \item 1909, EUA: \alert<1>{Junior Wireless Club}
    \item 1912, EUA: Promulgação do \alert<2>{Radio Act}
    \item 1914, EUA: Fundação da \alert<3>{ARRL}
    \item 1924, Brasil: \alert<4>{Regulamentação do rádio} e radioamadorismo, 5 de novembro
    \item 1925, Paris: fundação da \alert<5>{IARU}
    \item 1934: Fundação da \alert<6>{LABRE}
    \end{itemize}

    \column{.45\textwidth}
    \foto{antenna1}{Detalhe da capa do número 1, ano 1 da revista \emph{Antenna}, 1926}
  \end{columns}
\end{frame}

\begin{frame}{Radioamadorismo: um hobby organizado}
  O radioamadorismo, assim como o montanhismo organizado, é \alert{federado}:

  \medskip
  \centering
  \figura{organizacoes}
\end{frame}

%%% Local Variables:
%%% TeX-master: "../presentation"
%%% fill-column: 100
%%% End:
